\begin{activity}{Visualizing Magnetic Fields}

\section*{Purpose}

This activity develops physical intuition for magnetic fields as vector fields generated by dipoles. You will explore how field geometry depends on source orientation, strength, and proximity.

\section*{Materials}
\begin{itemize}
  \item Disc magnets
  \item Paper sheets
  \item Iron filings
\end{itemize}

\section*{Part A: Single disc magnet --- vertical orientation}

Place a disc magnet flat on the table and sprinkle iron filings over a sheet of paper on top.

\begin{enumerate}
  \item Describe the symmetry of the field pattern.
  \item Where are the field gradients strongest?
  \item If this magnet were buried at depth, where would you expect the strongest anomaly?
\end{enumerate}

\vspace{3cm}
\noindent\textbf{Sketch:}

\vspace{4cm}

\section*{Part B: Single disc magnet --- horizontal orientation}

Stand the same magnet on its edge and repeat the experiment.

\begin{enumerate}
  \item How does the field symmetry differ from Part A?
  \item Where are positive and negative lobes located?
  \item How would this differ from a gravity anomaly produced by the same body?
\end{enumerate}

\vspace{3cm}
\noindent\textbf{Sketch:}

\vspace{4cm}

\section*{Part C: Stacked disc magnets}

Stack several disc magnets and repeat Parts A and B.

\begin{enumerate}
  \item How does stacking affect field strength?
  \item Does stacking change anomaly shape or only amplitude?
  \item What geological parameter does stacking mimic?
\end{enumerate}

\section*{Reflection}

Why are magnetic field patterns sensitive to orientation, while gravity fields are not?

\end{activity}